\begin{trapbox}

	\begin{description}[style=nextline, leftmargin=2.8em, labelsep=0.5em, itemsep=0.35em]
		\item[\textbf{\textcolor{CTrap}{易错点一（实根前提）}}]
			不检验判别式 $\Delta$ 是否非负，直接给出两根之和与积。
		\item[\textbf{\textcolor{CTrap}{正确理解（验$\Delta$）：}}]
			使用韦达定理前必须先计算 $\Delta=b^2-4ac$，确保 $\Delta\geq0$（实根情形）。即使 $\Delta<0$ 公式形式仍成立，但高中试题常要求实根，故需先验证。
		\item[\textbf{\textcolor{CTrap}{错因分析（忽略条件）：}}]
			只记公式不注意适用条件，解题时盲目套用。
	\end{description}

	\medskip
	\begin{description}[style=nextline, leftmargin=2.8em, labelsep=0.5em, itemsep=0.35em]
		\item[\textbf{\textcolor{CTrap}{易错点二（和号遗漏）}}]
			将 $x_1+x_2$ 误记为 $\dfrac{b}{a}$，丢失负号。
		\item[\textbf{\textcolor{CTrap}{正确理解（符号辨识）：}}]
			正确公式为 $x_1+x_2=-\dfrac{b}{a}$，负号来源于求根公式中分子 $-b$。可记忆为“和是负比”。
		\item[\textbf{\textcolor{CTrap}{错因分析（记忆混淆）：}}]
			与一次项系数符号混淆，缺乏对公式来源的理解。
	\end{description}

	\medskip
	\begin{description}[style=nextline, leftmargin=2.8em, labelsep=0.5em, itemsep=0.35em]
		\item[\textbf{\textcolor{CTrap}{易错点三（积缺分母）}}]
			对于 $ax^2+bx+c=0$（$a\neq1$），错误认为 $x_1x_2=c$ 而忘记除以 $a$。
		\item[\textbf{\textcolor{CTrap}{正确理解（归一化处理）：}}]
			正确公式为 $x_1x_2=\dfrac{c}{a}$。若 $a=1$ 才为 $c$，否则必须除以 $a$。
		\item[\textbf{\textcolor{CTrap}{错因分析（轻视系数）：}}]
			只记忆了首项为 $1$ 时的特例，忽略一般形式。
	\end{description}

\end{trapbox}
